\documentclass[12pt,reqno]{amsart}
\usepackage{amsmath,amssymb,amsthm}
\usepackage{mathptmx}
\usepackage[scaled=0.90]{helvet}
\usepackage{courier}
\usepackage[T1]{fontenc}
\usepackage{enumitem}
\usepackage[hidelinks]{hyperref}
\pagestyle{plain}

\setlength{\oddsidemargin}{0cm}
\setlength{\evensidemargin}{0cm}
\setlength{\topmargin}{0in}
\setlength{\headheight}{0pt}
\setlength{\headsep}{0pt}
\setlength{\footskip}{.3in}
\setlength{\textheight}{9.2in}
\setlength{\textwidth}{6.5in}
\setlength{\parindent}{0pt}
\setlist[enumerate,1]{leftmargin=*,itemsep=1.1em}
\setlist[enumerate,2]{label=(\alph*),leftmargin=1.7em,itemsep=.45em,topsep=.4em}

\newcommand{\C}{\mathbb{C}}
\newcommand{\R}{\mathbb{R}}

\title{Problem Set 3 \\ MATH 411: Honors Complex Variables, Fall 2026}
\date{last compiled: \today}

% Coverage: material excluded from Midterm 1 through Monday, Oct. 26,
% 2026 (Lectures 13--20), according to the current course schedule.
% Suggested completion date: Monday, October 26, 2026.
% Sources (paths relative to the course workspace):
% 1: Old Homework/Homework/Fall2021-MATH411-Homework4.tex, 1(a)--(b).
% 2: same source, 2; adapted to the fixed-endpoint definition in Lecture 17.
% 3: same source, 3; f is required to be nonconstant to avoid a trivial
%    counterexample whose image is not open.
% 4: same source, 4--5; argument intervals replace ambiguous branch names.
% 5: same source, 6; sharpened to specified domains and a global obstruction.
% 6: Old Homework/Homework/Homework5_Fall2021.tex, 1.
% 7: Old Homework/Homework/Homework5_Fall2021.tex, 3--4.
% Residue and winding-number questions from later old sets are omitted.

\begin{document}
\maketitle

\begin{center}
Suggested completion date: Monday, October 26, 2026.
\end{center}

These are practice problems and are not to be turned in. Justify your
answers and all limiting arguments. Throughout, homotopies of paths
\emph{fix their endpoints}, and a simply connected open set is required
to be nonempty and connected.

\begin{enumerate}

\item \textbf{Real integrals by contour integration.}
Use Cauchy's theorem to evaluate the following improper integrals.
\begin{enumerate}
\item
\[
 \int_0^\infty \cos(t^2)\,dt.
\]
\emph{Hint:} Integrate $e^{-z^2}$ around the boundary of the sector
$0\leq\arg z\leq\pi/4$, $|z|\leq R$.
You may use
\[
 \int_{-\infty}^{\infty}e^{-x^2}\,dx=\sqrt{\pi},\qquad
 \sin u\geq\frac{2u}{\pi}\quad(0\leq u\leq\pi/2).
\]
Show that the integral over the circular arc tends to zero; a bound
using only the maximum of the integrand on the arc will not suffice.
\item For $a,b>0$, evaluate
\[
 \int_0^\infty e^{-ax}\cos(bx)\,dx.
\]
\emph{Hint:} Put $A=\sqrt{a^2+b^2}$ and choose
$0<\theta<\pi/2$ with $\cos\theta=a/A$ and $\sin\theta=b/A$.
Integrate $e^{-Az}$ around the sector of radius $R$ and angle $\theta$.
\end{enumerate}

\item \textbf{Star-shaped open sets.}
A nonempty open set $U\subseteq\C$ is \emph{star-shaped with respect
to $a\in U$} if the segment $[a,z]$ lies in $U$ for every $z\in U$.
\begin{enumerate}
\item Prove that $U$ is path connected.
\item For a closed continuous curve $\gamma:[0,1]\to U$ based at $a$,
construct a homotopy with fixed endpoints from $\gamma$ to the constant
curve at $a$. Deduce that $U$ is simply connected.
\end{enumerate}
For the last deduction, you may use the elementary fact that in a path
connected set it suffices to contract loops at one fixed base point.
A connecting path and its reverse transfer the contraction to any
other base point.

\item \textbf{Images of simply connected sets.}
Let $U$ be a simply connected open set and let $f$ be a nonconstant
holomorphic function on $U$. Must $f(U)$ be simply connected?
Investigate this using
\[
 H=\{z\in\C:\operatorname{Im}z>0\},\qquad f(z)=e^{2\pi i z}.
\]
Determine $f(H)$ exactly and justify your conclusion about its simple
connectivity, for example by integrating $1/z$ around a suitable circle.

\item \textbf{Two branches of the logarithm.}
Let $L_+$ be the logarithm on $\C\setminus[0,\infty)$ with argument
in $(0,2\pi)$, and let $L_-$ be the logarithm on
$\C\setminus(-\infty,0]$ with argument in $(-\pi,\pi)$.
For \emph{each} branch $L$, compute
\[
 L(-1+i),\qquad \exp\bigl(iL(-i)\bigr),\qquad L(-1-i).
\]
Here $(-i)^i$ is defined relative to the chosen branch as
$\exp(iL(-i))$. Explain how $L_+$ and $L_-$ differ on the upper and
lower half-planes.

\item \textbf{Holomorphic roots and the choice of domain.}
Fix an integer $n\geq2$.
\begin{enumerate}
\item Let $U\subseteq\C\setminus\{0\}$ be a simply connected open
set. Use a holomorphic logarithm on $U$ to construct a holomorphic
function $g$ satisfying $g(z)^n=z$ on $U$.
\item For $\theta\in\R$, give such a branch on the slit plane
\[
 U_\theta=\C\setminus\{re^{i\theta}:r\geq0\},
\]
specifying the interval used for the argument.
\item Show that there is no continuous function $g$ on
$\C\setminus\{0\}$ satisfying $g(z)^n=z$ everywhere.
\emph{Hint:} On the unit circle, consider
$g(e^{it})e^{-it/n}$ for $0\leq t\leq2\pi$.
\end{enumerate}

\item \textbf{Polynomial growth and Cauchy estimates.}
Let $f$ be entire and put
\[
 M_f(R)=\max_{|z|=R}|f(z)|.
\]
Suppose there are a constant $C>0$, an integer $k\geq0$, and radii
$R_j\to\infty$ such that $M_f(R_j)\leq C R_j^k$ for every $j$.
Use Cauchy's estimates for the Taylor coefficients at $0$ to prove
that $f$ is a polynomial of degree at most $k$ (including the zero
polynomial). What does the case $k=0$ say?

\item \textbf{Laurent expansions.}
\begin{enumerate}
\item Find a series in integer powers of $z$ representing
\[
 f(z)=\frac{1}{(z-1)(z-2)}
\]
in each of the regions $|z|<1$, $1<|z|<2$, and $|z|>2$.
Justify convergence in each region.
\item Find the first four nonzero terms of the Laurent expansion of
\[
 \frac{e^z}{z(1+z^2)}
\]
valid for $0<|z|<1$.
\end{enumerate}

\end{enumerate}
\end{document}
